\documentclass[ctcc_full_documentation.tex]{subfiles}
\begin{document}
\chapter{Congestion and Curtailment Benefits Calculation Documentation}

\section{Introduction}

The \texttt{congestion\_curtailment\_reduction.py} script calculates the benefits from reducing transmission congestion and renewable energy curtailment. It allocates effective capacity relief between curtailment and congestion to avoid double-counting, then monetizes both benefits over the project lifetime.

\section{Method Overview}

\subsection{Purpose}

This script calculates congestion and curtailment benefits by:
\begin{enumerate}
    \item Determining effective capacity relief (greenfield vs. reconductoring)
    \item Allocating capacity relief between curtailment and congestion
    \item Calculating congestion reduction energy and costs
    \item Calculating curtailment reduction energy and benefits
    \item Applying saturation factors (haircuts) to account for diminishing returns
    \item Computing residual exceedance costs (unmet capacity needs; one cost from congestion and/or curtailment)
    \item Calculating present values of all benefits and costs
\end{enumerate}

\subsection{Calculation Approach}

The calculation follows this structure:

\textbf{Effective Capacity Relief:}
\begin{align}
\Delta C_{eff} &= \begin{cases}
\phi \times \text{Capacity} & \text{(Greenfield)} \\
\text{Capacity} - \text{Old Capacity} & \text{(Reconductoring)}
\end{cases}
\end{align}

where $\phi$ is the flow factor.

\textbf{Curtailment Benefits:}
\begin{align}
E_{curt} &= H_c \times \min(\Delta C_{eff}, C_c) \\
\text{Benefit}_{curt} &= E_{curt} \times \lambda_c
\end{align}

where $H_c$ is curtailment hours, $C_c$ is average curtailment MW, and $\lambda_c$ is curtailment price.

\textbf{Congestion Benefits:}
\begin{align}
E_{cong} &= E_{cong,bc} + E_{cong,non} + E_{near} \\
\text{Benefit}_{cong} &= E_{cong} \times \lambda_{cong} \times (1 - \text{saturation})
\end{align}

where the energy is split between overlap hours (with curtailment) and non-overlap binding hours.

\textbf{Residual Exceedance:}
Residual exceedance is one cost that can have two sources---congestion binding hours and/or curtailment hours---so it works for congestion-only, curtailment-only, or both. When only one of congestion or curtailment is used, the other term is zero; when both are used, the user should set \texttt{residual\_exceedance\_value} if the single default (\$/MWh) is not appropriate.

\begin{align}
E_{res,cong} &= (H_{bc} + H_{bnon}) \times \max(0, X_{cong} - \Delta C_{eff}) \\
E_{res,curt} &= H_c \times \max(0, C_c - \Delta C_{eff}) \\
E_{residual} &= E_{res,cong} + E_{res,curt} \\
\text{Cost}_{residual} &= E_{residual} \times \lambda, \quad \lambda = \text{\texttt{residual\_exceedance\_value}} \text{ if set, else } \lambda_{cong}
\end{align}

\subsection{Inputs and Outputs}

\textbf{Inputs:}
\begin{itemize}
    \item Project technical details (capacity, old capacity, reconductoring status)
    \item Congestion parameters (binding hours, average exceedance, flow factor, saturation factor, congestion price; \texttt{residual\_exceedance\_value} optional, null = use average congestion price)
    \item Curtailment parameters (curtailment hours, average curtailment MW, curtailment price, saturation factor)
    \item Financing parameters (WACC, timeline)
\end{itemize}

\textbf{Outputs:}
\begin{itemize}
    \item Annual and lifetime congestion benefits (raw and haircut)
    \item Annual and lifetime curtailment benefits (raw and haircut)
    \item Residual exceedance costs
    \item Congestion and curtailment costs during delay/construction (allocated, no double count)
    \item Present values of all metrics
\end{itemize}

\subsection{Integration with Other Scripts}

This script integrates with:
\begin{itemize}
    \item \texttt{smart\_loaders.py} --- \texttt{load\_congestion\_curtailment\_reductions}, \texttt{load\_project\_technical\_details}, \texttt{load\_financing\_details}, \texttt{get\_project\_data\_raw}
    \item \texttt{yaml\_loaders.py} --- \texttt{ProjectTechnicalDetails} (via centralized loader)
    \item \texttt{financial\_utils.py} --- \texttt{calculate\_present\_value}, \texttt{calculate\_cod\_year}
    \item \texttt{constants.py} --- \texttt{MIN\_DISCOUNT\_RATE}
    \item \texttt{smart\_output.py} --- \texttt{CTCCOutputManager} (\texttt{add\_congestion\_curtailment}, \texttt{write\_batch\_summary})
\end{itemize}

\section{Key Functions}

\subsection{\texttt{allocate\_curtailment\_then\_congestion(DC\_eff\_mw, ...)}}

Allocates effective capacity relief between curtailment and congestion to avoid double-counting.

\textbf{Parameters:}
\begin{itemize}
    \item \texttt{DC\_eff\_mw} (float): Effective capacity relief in MW
    \item \texttt{binding\_hours\_total} (float): Total binding constraint hours
    \item \texttt{average\_exceedance\_mw} (float): Average exceedance in MW
    \item \texttt{curtailment\_hours\_total} (float): Total curtailment hours
    \item \texttt{average\_curtailment\_mw} (float): Average curtailment in MW
    \item \texttt{average\_curtailment\_price} (float): Average curtailment price
\end{itemize}

\textbf{Returns:}
\begin{itemize}
    \item Dictionary with curtailment benefits, overlap metrics, and remaining capacity for congestion
\end{itemize}

\textbf{Key Logic:}
\begin{lstlisting}
Hb = max(0.0, binding_hours_total)
Hc = max(0.0, curtailment_hours_total)
# Curtailment relief (MWh/yr)
E_curt = Hc * min(DeltaC_eff_mw, Cc)
curt_benefit = E_curt * lambda_c

if Hb <= 0:
    return {..., "H_bc": 0, "H_bnon": 0, "DeltaC_remain_bc_mw": DeltaC_eff_mw}

# Overlap proxy between binding and curtailment time
theta = min(1.0, Hc / Hb) if Hc > 0 else 0.0
H_bc = theta * Hb
H_bnon = (1.0 - theta) * Hb
DeltaC_used_bc = min(DeltaC_eff_mw, Cc)
DeltaC_remain_bc = max(0.0, DeltaC_eff_mw - DeltaC_used_bc)
# Returns E_curt_mwh_yr, curtailment_benefit_$_yr, theta_overlap, H_bc, H_bnon, DeltaC_remain_bc_mw
\end{lstlisting}

\subsection{\texttt{calculate\_congestion\_reduction\_costs(...)}}

Main function that calculates congestion and curtailment benefits.

\textbf{Parameters:}
\begin{itemize}
    \item Multiple parameters including reconductoring status, capacities, congestion/curtailment parameters, financing parameters
\end{itemize}

\textbf{Returns:}
\begin{itemize}
    \item Dictionary with all calculated benefits, costs, and present values
\end{itemize}

\textbf{Key Logic:}
\begin{lstlisting}
if not reconductoring:
    effective_capacity_relief = max(0, flow_factor * capacity_mw)
else:
    effective_capacity_relief = max(0, capacity_mw - old_capacity_mw)

alloc = allocate_curtailment_then_congestion(DeltaC_eff_mw=effective_capacity_relief, ...)
DeltaC_rem = alloc["DeltaC_remain_bc_mw"]
H_bc, H_bnon = alloc["H_bc"], alloc["H_bnon"]

E_cong_bc = H_bc * min(DeltaC_rem, average_exceedance)
E_cong_non = H_bnon * min(effective_capacity_relief, average_exceedance)
E_near = near_binding_hours * min(effective_capacity_relief, near_average_exceedance)
energy_congestion_reduction = E_cong_bc + E_cong_non + E_near

annual_congestion_reduction_cost_raw = energy_congestion_reduction * average_congestion_price
annual_congestion_reduction_cost_haircut = (1 - saturation_factor) * annual_congestion_reduction_cost_raw

total_binding_hours = H_bc + H_bnon
energy_residual_congestion = total_binding_hours * max(0, average_exceedance - effective_capacity_relief)
energy_residual_curtailment = curtailment_hours_total * max(0, average_curtailment_mw - effective_capacity_relief)
energy_residual_exceedance = energy_residual_congestion + energy_residual_curtailment
residual_exceedance_value_used = residual_exceedance_value or average_congestion_price
annual_residual_exceedance_cost = energy_residual_exceedance * residual_exceedance_value_used
\end{lstlisting}

\subsection{\texttt{main()}}

Main execution function that orchestrates the calculation and output.

\section{Example}

The following example demonstrates a typical usage:

\begin{lstlisting}
# Example: 500 MW Greenfield Overhead AC line
# 
# Inputs (from YAML):
#   - Capacity: 500 MW
#   - Flow factor: 0.7
#   - Binding hours: 2000 hours/year
#   - Average exceedance: 300 MW
#   - Curtailment hours: 1500 hours/year
#   - Average curtailment: 200 MW
#   - Congestion price: $50/MWh
#   - Curtailment price: $30/MWh
#   - Saturation factor: 0.2 (20% haircut)
#   - Project lifetime: 50 years
#
# Calculation:
#   Effective capacity relief = 0.7 * 500 = 350 MW
#   
#   Curtailment:
#     E_curt = 1500 * min(350, 200) = 300,000 MWh/year
#     Benefit = 300,000 * $30 = $9M/year
#   
#   Congestion (after curtailment allocation):
#     Overlap theta = min(1.0, 1500/2000) = 0.75
#     H_bc = 0.75 * 2000 = 1500 hours
#     H_bnon = 0.25 * 2000 = 500 hours
#     DC_remain = 350 - 200 = 150 MW
#     E_cong_bc = 1500 * min(150, 300) = 225,000 MWh/year
#     E_cong_non = 500 * min(350, 300) = 150,000 MWh/year
#     E_cong = 375,000 MWh/year
#     Benefit_raw = 375,000 * $50 = $18.75M/year
#     Benefit_haircut = (1 - 0.2) * $18.75M = $15M/year
#   
#   Residual exceedance (congestion contribution here; curtailment-only: E_res,cong=0,
#     E_res,curt = H_c * max(0, C_c - Delta_C_eff)):
#     E_residual = 1500 * max(0, 300-150) + 500 * max(0, 300-350) 
#                = 225,000 MWh/year
#     Cost = 225,000 * $50 = $11.25M/year
#
# Output:
#   - Annual curtailment benefit: $9M
#   - Annual congestion benefit (haircut): $15M
#   - Annual residual exceedance cost: $11.25M
#   - Net annual benefit: $12.75M
\end{lstlisting}

\textbf{Integration Example:}

\begin{lstlisting}
# When called from ctcc.py (main() builds results from CongestionReductionResults):
csv_manager = CTCCOutputManager()
results = {
    "congestion_benefit_annual": congestion_results.annual_congestion_reduction_cost_raw,
    "congestion_benefit_nominal": congestion_results.lifetime_congestion_reduction_cost,
    "congestion_benefit_pv": congestion_results.lifetime_congestion_reduction_cost_pv,
    "curtailment_benefit_annual": congestion_results.annual_curtailment_benefit,
    "curtailment_benefit_pv": congestion_results.lifetime_curtailment_benefit_pv,
    "residual_exceedance_annual": congestion_results.annual_residual_exceedance_cost,
    "residual_exceedance_pv": congestion_results.lifetime_residual_exceedance_cost_pv,
    # ... delay costs, haircut variants, physical metrics
}
csv_manager.add_congestion_curtailment(results)
csv_manager.write_batch_summary()
\end{lstlisting}

\section{Key Implementation Details}

\subsection{Double-Counting Avoidance}

The script prevents double-counting by allocating capacity relief first to curtailment, then to congestion. The overlap factor (theta) accounts for hours when both curtailment and congestion occur. This same allocation is used for delay-period opportunity costs.

\subsection{Saturation Factors}

Saturation factors (haircuts) account for diminishing returns as more capacity is added. They reduce the monetized benefits to reflect that not all theoretical benefits are realized.

\subsection{Residual Exceedance}

The script calculates residual exceedance costs, representing unmet capacity when effective relief is less than the constraint (exceedance or curtailment). Residual exceedance can come from (1) congestion binding hours, (2) curtailment hours, or both. One \$/MWh price is applied (default: average congestion price; the user can set \texttt{residual\_exceedance\_value} when both apply). This is important for partial solutions (e.g., reconductoring or a line that does not fully relieve curtailment). In CTCC, residual exceedance is a system/societal cost categorized under energy/emissions (not under the utility or ratepayer perspective).

\subsection{Flow Factor}

For greenfield projects, the flow factor ($\phi$) represents the fraction of nameplate capacity that effectively relieves congestion, accounting for network constraints and power flow physics.

\subsection{Near-Binding Hours}

The script includes near-binding hours in a way that is methodologically consistent with binding hours: relief is capped by \texttt{near\_average\_exceedance}. The formula is $E_{near} = H_n \times \min(\Delta C_{eff}, X_n)$ where $H_n$ is near-binding hours and $X_n$ is near-average exceedance (MW).

\subsection{Delay and Construction Period Costs}

The script calculates delay-period opportunity costs using the same curtailment-first allocation as operations, so overlapping hours are not double counted. This represents the value of forgone relief while the project is delayed. Present values use \texttt{financial\_utils.calculate\_present\_value} with \texttt{start\_year=calculate\_cod\_year(delay\_years, construction\_years)} for benefits (annuity starts at COD); delay/construction costs use \texttt{start\_year=1} over \texttt{delay\_years + construction\_years}.

\end{document}

